\documentclass[10pt,a4paper]{amsart}
\usepackage[margin=19mm]{geometry}
\usepackage{amsmath,amssymb,mathtools}
\usepackage[scr=rsfs,cal=euler]{mathalfa}
\usepackage{xfrac}
\usepackage[unicode,hidelinks,bookmarks=false]{hyperref}
\newcommand{\ie}{i.e.\ }
\newcommand{\cf}{cf.\ }
\newcommand{\resp}{respectively\ }
\newcommand{\Subsection}{\subsection}
\providecommand{\nobreakdash}{\nobreak\hbox{-}}
\setlength{\emergencystretch}{2em}
\multlinegap0pt
\usepackage{amssymb}
\usepackage{bm}
\usepackage[shortlabels]{enumitem}
\newtheorem{theorem}{Theorem}
\newcommand{\nn}{\nonumber}
\title{Laurent sequences, extended Rota Algebras \\ and Categorical Discretization \\ of Dynamical Systems}
\author{Miguel A. Rodr\'iguez, Piergiulio Tempesta}
\date{Source: arXiv:2510.15489v3 (CC BY 4.0)}
\begin{document}
\maketitle

\noindent\emph{Source and licence.} Laurent sequences, extended Rota Algebras \\ and Categorical Discretization \\ of Dynamical Systems. Version arXiv:2510.15489v3 (CC BY 4.0), distributed by arXiv under the Creative Commons Attribution 4.0 International licence, http://creativecommons.org/licenses/by/4.0/, as recorded in the arXiv licence metadata for this exact version and shown on the abstract page https://arxiv.org/abs/2510.15489v3. Full licence notice: \texttt{licenses/arxiv-2510.15489-CC-BY-4.0.txt}.\par\smallskip

\noindent\textit{Editorial scope.} Counted unit: the article's theorem main4 (source0.tex lines 1686--1722). The article's complete original proof of it is delivered with it (source0.tex lines 1723--1826, one \texttt{proof} environment of 104 lines). No mathematics is written, repaired or re-derived: the statement and the proof are the article's own lines, in the article's own notation, and no retained line is substituted or re-flowed.  Retained uncounted as the source-local context the proof needs: lines 1569--1573.  Nothing was omitted from any retained span, so the package carries no omission record.  No retained line contains a citation, so the wrapper carries no bibliography and no bibliography entry.  No retained line contains a figure environment, an image-inclusion command or drawing code, so no image is delivered and the package declares no asset dependency.  Editorial formatting only, in addition: an AMS \texttt{amsart} preamble that restates the article's own declarations and macros used by the retained lines, namely the environments \texttt{theorem} on separate counters, together with the standard packages those lines need.  The retained environments print in this document as Theorem 1 in order of appearance, whereas the article prints the same items with the numbering of its own full document.  The complete native source of the article as distributed in the arXiv e-print for this exact version is bundled unchanged as \texttt{sources/187261/source0.tex}.
\begin{equation}
\label{eq:main3}
lin(\partial , y):=a_{N}(x) \frac{d^{N}}{dx^{N}}y + a_{N-1}(x)
\frac{d^{N-1}}{dx^{N-1}}y +\ldots +a_{1}(x) \frac{d}{dx} y+a_{0}(x) y +g(x)=0,
\end{equation}

\begin{theorem}
\label{main4}
Consider an ordinary differential equation of order
$m\in \mathbb{N}\setminus \{0\}$ of the form
\begin{equation}
\label{eq:main4}
eq(\partial , y):=\frac{d^{m}}{dx^{m}}y-\sum _{r=0}^{N} c_{r}(x) y^{r}=0,
\end{equation}
where $N, m\in \mathbb{N}\setminus \{0\}$, $c_{r}(x)$,
$r=0,\ldots ,N$ are real functions admitting at most polar singularities at
$x=0$ of the form:
\begin{equation}
c_{r}(x)=\sum _{j=-\ell _{r}}^{\rho _{r}}\beta _{rj} x^{j}.
\label{eq8.12}
\end{equation}
Assume that eq.~\eqref{eq:main4} has a solution possessing a polar singularity
of order $s$ at the origin:
\begin{equation}
\label{eq:npolsol}
y(x)= \sum _{k=1}^{s}\frac{\zeta _{-k}}{x^{k}},
\end{equation}
where
$s\geq \max \{\ell _{0}+\rho _{0},\ldots ,\ell _{N}+\rho _{N} \}$. Then,
the difference equation
\begin{equation}
\label{eq:8.14}
\sum _{i=0}^{m}(-1)^{m-i}\binom{m}{i}u_{n+i}=\sum _{r=0}^{N} \sum _{j=-
\ell _{r}}^{\rho _{r}} \sum _{k_{1},\cdots ,k_{r}=-\infty}^{-1}
\beta _{rj} \, \zeta _{k_{1}}\cdots \zeta _{k_{r}}p_{K_r+j}(n)
\end{equation}
with $K_{r}=\sum _{i=1}^{r} k_{i}$, defined over a uniform lattice in the positive real line, has a
solution of the form
\begin{equation}
\label{eq:8.15}
u_{n}=\sum _{k=1}^{s}\frac{(n-1)!}{(n+k-1)!}\zeta _{-k}.
\end{equation}
\end{theorem}

\begin{proof}
\label{eq:8.23}
The discretization runs as follows. The function $y^{r}(x)$, is written
as a Laurent expansion centred at the origin. Let us write formally
\begin{equation}
y^{r}\to \left (\sum _{k_{1}=-\infty}^{-1}\zeta _{k_{1}} p_{k_{1}}(n)
\right )*\cdots * \left (\sum _{k_{r}=-\infty}^{-1}\zeta _{k_{r}} p_{k_{r}}(n)
\right )= \sum _{k_{1},\cdots ,k_{r}=-\infty}^{-1}\zeta _{k_{1}}
\cdots \zeta _{k_{r}}p_{K_{r}}(n),
\label{eq8.16}
\end{equation}
where $K_{r}=\sum _{i=1}^{r} k_{i}$, and $\zeta _{k_{i}}=0$ for
$k_{i}=-s-1,\ldots , -\infty $. Thus,
\begin{eqnarray}
c_{r}(x) y^{r} &\to & \sum _{j=-\ell _{r}}^{\rho _{r}}\beta _{rj} p_{j}(n)*
\sum _{k_{1},\cdots ,k_{r}=-\infty}^{-1}\zeta _{k_{1}}\cdots \zeta _{k_{r}}p_{K_{r}}(n)
=
\label{eq8.17}
\\
\nonumber
&=&\sum _{j=-\ell _{r}}^{\rho _{r}} \sum _{k_{1},\cdots ,k_{r}=-
\infty}^{-1} \beta _{rj} \, \zeta _{k_{1}}\cdots \zeta _{k_{r}}p_{K_{r}+j}(n).
\end{eqnarray}
Then, the discretized function, $u_{n}$, is written as
\begin{equation}
u_{n}=\sum _{k=1}^{s}\frac{\zeta _{-k}}{p^{-}_{k}(n)}.
\label{eq8.18}
\end{equation}
The derivatives $y^{(m)}$, $m=1,\ldots , r$, are discretized according
to the operator $\Delta $:
\begin{equation}
\Delta ^{m} u_{n}=\sum _{i=0}^{m}(-1)^{i}\binom{m}{i}u_{n+i}.
\label{eq8.19}
\end{equation}
The constants $\zeta _{-k}$ satisfy a linear algebraic system:
\begin{equation}
\label{eq:syst}
\begin{cases}
u_{n+1} = \sum _{k=1}^{s}\frac{\zeta _{-k}}{p^{-}_{k}(n+1)}= \sum _{k=1}^{s}
A_{n+1,k}\,\zeta _{-k},
\\
\hspace{3mm}
\vdots
\\
u_{n+s} = \sum _{k=1}^{s}\frac{\zeta _{-k}}{p^{-}_{k}(n+s)}= \sum _{k=1}^{s}
A_{n+s,k}\,\zeta _{-k},
\end{cases}
\end{equation}
where $A_{j k}=\frac{(j-1)!}{(j+k-1)!}$, $k=1\ldots ,s$,
$j=n+1\ldots ,n+s$. Unlike in the previous theorems, where the corresponding
coefficient matrix acquires a triangular form, in the present case the
solution of system  \eqref{eq:syst} cannot be easily written in a closed
form in terms of the variables $u_{n+i}$. We obtain
\begin{equation}
\zeta _{-k}= \sum _{\lambda =1}^{s} \big(A^{-1}\big)_{k,\lambda}u_{n+
\lambda}.
\label{eq8.21}
\end{equation}
Thus, the discrete equation reads:
\begin{eqnarray}\nn
\sum _{i=0}^{m}(-1)^{i}\binom{m}{i}u_{n+i} = \sum _{r=0}^{N} \sum _{j=-
\ell _{r}}^{\rho _{r}} \sum _{k_{1},\cdots ,k_{r}=-\infty}^{-1} \sum _{
\lambda _{1},\ldots ,\lambda _{r}=1}^{s} \beta _{rj} \big(A^{-1}\big)_{k_{1},
\lambda _{1}}\cdots \big(A^{-1}\big)_{k_{r},\lambda _{r}} \cdot
\label{eq8.22}
\\
%\nonumber
\cdot \,\, u_{n+\lambda _{1}}\cdots u_{n+\lambda _{r}}p_{K_{r}+j}(n),
\end{eqnarray}
where
\begin{equation}
p_{K_r+j}(n)=
\begin{cases}
\dfrac{n!}{(n-K_r-j)!}, & K_r+j\geq 0,
\\[10pt]
\dfrac{(n-1)!}{(n-K_r-j-1)!}, & K_r+j <0.
\end{cases}
\label{eq8.23}
\end{equation}

Let us prove that the series  \eqref{eq:8.15} is solution of eq.~\eqref{eq:8.14}. To this aim, we observe that the morphism
$\rho _{\partial ,\Delta}$ carries on
\begin{equation}
\sum _{k=1}^{s} \zeta _{-k} x^{-k} \longrightarrow \sum _{k=1}^{s}
\zeta _{-k} \frac{1}{p_{k}^{-}(n)}.
\label{eq8.24}
\end{equation}
Thus, the action of the functor $F$ carries any solution $y$ of eq.~\eqref{eq:main3},
defined on the algebra $(\mathcal{H},+,\cdot )$, into a solution
$u_{n}$ of the corresponding equation
\begin{equation*}
eq(\Delta ,u,*_{\Delta})=\nu _{\partial ,\Delta}(eq(\partial , y,
\cdot )),
\end{equation*}
defined on the extended Rota algebra $(\mathcal{H},+,*_{\Delta})$. According
to the previous analysis, we represent this equation on the lattice
$\mathcal{L}$. Then, any finite Laurent series solution  \eqref{eq:npolsol} of  eq.~\eqref{eq:main4} is carried into a solution of
the form  \eqref{eq:8.15} of the equation
\begin{equation*}
eq(\mathcal{Q},u,*_{\mathcal{Q}})=\nu _{\partial ,\Delta}(eq(
\partial ,y,\cdot )).
\end{equation*}
This equation coincides with eq.~\eqref{eq:8.14}.
\end{proof}

\end{document}
